\documentclass[12pt,a4paper,titlepage,oneside]{article}
\usepackage{isabelle,isabellesym}

% further packages required for unusual symbols (see also
% isabellesym.sty), use only when needed

%\usepackage{amssymb}
  %for \<leadsto>, \<box>, \<diamond>, \<sqsupset>, \<mho>, \<Join>,
  %\<lhd>, \<lesssim>, \<greatersim>, \<lessapprox>, \<greaterapprox>,
  %\<triangleq>, \<yen>, \<lozenge>

%\usepackage{eurosym}
  %for \<euro>

%\usepackage[only,bigsqcap]{stmaryrd}
  %for \<Sqinter>

%\usepackage{eufrak}
  %for \<AA> ... \<ZZ>, \<aa> ... \<zz> (also included in amssymb)

%\usepackage{textcomp}
  %for \<onequarter>, \<onehalf>, \<threequarters>, \<degree>, \<cent>,
  %\<currency>

\usepackage[utf8]{inputenc}
\usepackage[T1]{fontenc}
\usepackage{amsmath}
\usepackage{array}
\usepackage{amsfonts}
\usepackage{tikz}
\usepackage[inline]{enumitem}
\usepackage{varwidth}
\usepackage{lipsum} % just to automatically generate text
\usepackage[nottoc,numbib]{tocbibind}
\usepackage{mathpartir}

\usetikzlibrary{arrows}

% this should be the last package used
\usepackage{pdfsetup}

% urls in roman style, theory text in math-similar italics
\urlstyle{rm}
\isabellestyle{it}

% for uniform font size
%\renewcommand{\isastyle}{\isastyleminor}
%
\hyphenation{Isa-belle}

\newcommand{\HRule}{\rule{\linewidth}{0.5mm}}

\makeatletter
\newcommand\ackname{Acknowledgements}
\if@titlepage
  \newenvironment{acknowledgements}{%
      \titlepage
      \null\vfil
      \@beginparpenalty\@lowpenalty
      \begin{center}%
        \bfseries \ackname
        \@endparpenalty\@M
      \end{center}}%
     {\par\vfil\null\endtitlepage}
\else
  \newenvironment{acknowledgements}{%
      \if@twocolumn
        \section*{\abstractname}%
      \else
        \small
        \begin{center}%
          {\bfseries \ackname\vspace{-.5em}\vspace{\z@}}%
        \end{center}%
        \quotation
      \fi}
      {\if@twocolumn\else\endquotation\fi}
\fi
\makeatother

\begin{document}
\pagestyle{empty}

% \title{Formalizing \emph{Types and Programming Languages} in Isabelle/HOL}
% \author{Martin Desharnais}

\begin{titlepage}
  \begin{center}
    % Upper part of the page. The '~' is needed because \\
    % only works if a paragraph has started.
    \includegraphics[width=6cm]{./logo-ETS.png}~\\[1cm]

    \textsc{\LARGE École de technologie supérieure}\\[1.5cm]
    \textsc{\Large Bachelor's Thesis}\\[0.5cm]

    % Title
    \HRule \\[0.4cm]
    { \huge \bfseries Formalizing \emph{Types and Programming Languages} in Isabelle/HOL \\[0.4cm] }

    \HRule \\[1.5cm]

    % Author and supervisor
    \begin{center}
      \begin{tabular}{lr}
        Author:    & Martin Desharnais \\
        Professor: & David Labbé \\
        Advisor:   & Jasmin C. Blanchette \\
        Advisor:   & Dmitriy Traytel
      \end{tabular}
    \end{center}

    \vfill

    % Bottom of the page
    {\large \today}

  \end{center}
\end{titlepage}

\newpage
\null
\newpage
\pagenumbering{roman}
%\maketitle

\begin{abstract}
We formalize, using Isabelle/HOL, some languages present in the first two sections, namely ``Untyped
Systems'' and ``Simple Types'', of the book \emph{Types and Programming Languages} by
Benjamin~C.~Pierce. We first begin with a short tour of the $\lambda$-calculus, type systems and the
Isabelle/HOL theorem prover before attacking the formalization \emph{per se}. Starting with an
arithmetic expressions language offering booleans and natural numbers, we pursue, after a brief
digression to the "de Bruijn indices", to the untyped $\lambda$-calculus. We then return to a typed
variant of the arithmetic expression language before to conclude with the simply typed
$\lambda$-calculus.
\end{abstract}

\newpage
\null
\newpage

\begin{acknowledgements}
  First, I would like to thank Jasmin Blanchette, one of my advisors, who gave me the opportunity to
  move to Munich for an eight month intership as a research assistant. This experience gave me a
  taste of academic research and introduced me to the world of formal verifications and theorem
  provers. I benefited greatly from the numerous discussions we had and the advices he gave me.

  I am equally grateful to Dmitriy Traytel, my second advisor, for all the help he gave me for
  both my internship and my thesis. He taught me a lot and his help was invaluable for me to
  overcome some of the challenges I encountered.

  I want to thank all members of the research chair for the warm welcome the gave me. A special
  thought goes to Lars Noschinski, my officemate, for all the conversations we had and all the
  insights he gave me of the German culture.
\end{acknowledgements}

\newpage
\null
\newpage

\tableofcontents

\newpage
\null
\newpage

\pagenumbering{arabic}
\pagestyle{headings}

\section{Introduction}
% in the Isabelle/HOL theorem prover

This Bachelor's thesis deals with the formalisation of parts of the book \emph{Types and Programming
Languages} \cite{pierce-2002-TAPL}, hereafter abbreviated \emph{TAPL}, by Benjamin~C.~Pierce. The
formalization is performed using the Isabelle/HOL theorem prover. This work concentrates on four
languages, ranging from simple arithmetic expressions to the fully fledged $\lambda$-calculus,
present in the first two sections, namely ``Untyped Systems'' and ``Simple Types''.

The main motivation to have chosen this subject is the intersection of personal interest and of
opportunities provided by my internship at the chair for logic and verification at Technische
Universität München. Having gradually developed an interest for programming languages in the last
years, I was eager to learn more about the theory behind type systems. Pierce's book stood out as
a reference recommended for a deep introduction to the main elements of this field. Also, as part of
my internship, I worked on the implementation of the (co)datatype module in the Isabelle/HOL
theorem prover. Having experienced the implementer role, I also wanted to learn about the user
role and the process of formalization. Thus, the choice of this subject for this thesis was an
opportunity to fulfill both goals.

Before entering into the realm of formalizations, we first introduce the required background
(Section \ref{sec:background}) for $\lambda$-calculus, type systems and Isabelle/HOL. The
$\lambda$-calculus is a core calculus that captures the essential features of functional programming
languages.  That is, there exists a way to encode high level features such as recursion, datatypes,
records, etc. Such calculus can come, as programming languages do, in two variants: typed and
untyped. A type system is a syntactic method for automatically checking the absence of certain
erroneous behaviors. To formalize these, we used Isabelle/HOL, an interactive theorem prover based
on higher order logic. It resembles a functional programming language in that one can define
datatypes and functions. The difference is that it allows to postulate properties of the formerly
defined elements and to provide machine checked poofs that those properties are theorems.

The formalizations we perform all have a direct correspondence with chapters from \emph{TAPL}
(Section \ref{sec:structure-of-formalization}). We provide one Isabelle/HOL theory file per chapter.

The untyped arithmetic expressions language (Section \ref{sec:untyped-arith-expr}) serves as a
warm-up to experiment with the general structure of formalizations. It consists of Boolean
expressions and natural numbers. This simplicity allows us to concentrate on the translation to
Isabelle/HOL of the definitions found in the book and to accustom ourselves with the notation. Most
of our definitions and theorems closely follows the ones from the book. The main exceptions are
that we expose some hypotheses that are implicit in the book and that we use a different definition
of the mutli-step evaluation relation.

The formalization (Section \ref{sec:nameless-rep-of-terms}) of the nameless representation of terms,
also known as de Bruijn indices'', was not initially planned but arose from the need to use a
concrete representation for variables in the $\lambda$-calculus. Our formalization closely follows
the book.

While the previous formalizations are either a warm-up or a representation necessity, the untyped
$\lambda$-calculus (Section \ref{sec:untyped-lambda-calculus}) is the first core calculus we
formalized. We differ from the book in the definition of the evaluation relation. In \emph{TAPL},
the evaluation relation assumes that name clashes in variables are automatically solved by renaming
them and, thus, ignore this possibility from there on. Such an assumption is not accepted by
computer-verified proof. We use de Bruijn indices'' as representation for variables to encode this
assumption. Also, since the chapter in the book is more focused on explaining the
$\lambda$-calculus, it contains no meaningful theorems. Nevertheless, we revisit the properties
introduced with the arithmetic expressions language and either prove that they are still theorems or
disprove them.

The typed arithmetic expressions language (Section \ref{sec:typed-arith-expr}) is again a warm-up;
this time, to experiment with the formalization of a type system. Our formalization closely follows
the book.

The simply typed $\lambda$-calculus (Section \ref{sec:simply-typed-lambda-calculus}) is the second
core calculus we formalize. Here, we differ significantly from the book, mainly because of our use
of de Bruijn indices'' but also because of our representation of the typing context, we need to
adapt some lemmas and replace others. This is certainly the most challenging part of the
formalization, since we cannot follow the described proofs anymore and must find the right
assumptions for our lemmas. In spite of these differences, our formalization still respects the
spirit of the book since only the helper lemmas change, and the important theorems remain the same.

All sections combined, the formalization consists of 800 lines of definitions, theorems and
exercises proposed in the book. It is publicly
available\footnote{\url{https://github.com/authchir/log792-type-systems-formalization}} and can be
executed with Isabelle 2014.\footnote{\url{http://isabelle.in.tum.de/website-Isabelle2014/}} In this
report, we focus on the definitions and how the theorems are expressed. When relevant, we present
both the definitions from the book and our translation, highlighting and motivating the differences.
Some proofs are presented but not explained. For a deeper insight into the proofs, the best
methodology is to study the theory files in Isabelle.

% $ wc formalization.txt
%  1606  9203 62268 formalization.txt
% $ wc formalization.txt
%   806  4681 33955 formalization.txt

\section{Background}
\label{sec:background}

\subsection{Lambda-Calculus}
\label{sec:background-lambda-calculus}

The $\lambda$-calculus is a minimalist language, where everything is a functions, that can be use as
a core calculus capturing the essential features of complex programming languages. It was formulated
by Alonzo Church \cite{???} in the 1930s as a model of computability. At the same time, Alan Turing
was formulating what is now known as a Turing machine \cite{???} for the same purpose. It was later
proved that both systems are equivalent in expressive power \cite{???}.

As a programming language, it can be intriguing at first because everything reduces to function
abstraction and application. The syntax comprises three sorts of terms: variables, function
abstractions over a variable and applications of a term to an other. Those three constructs are
summarized in the following grammar of $\lambda$-terms using a BNF-like notation:
\begin{align*}
  t ::= & \\
    & x && \text{variable} \\
    & \lambda x. t && \text{abstraction} \\
    & t_1 \text{ } t_2 && \text{application}
\end{align*}

Following are a few standard $\lambda$-term shown as examples of how the grammar is actually
used:\footnote{To reduce the need for parenthesis, we use the following standard conventions:
\begin{enumerate*}[label=(\arabic*)]
  \item the body of a $\lambda$-abstraction expands as much as possible to the right and
  \item function application is left associative.
\end{enumerate*}}
\begin{align*}
  & \lambda x. x && \text{identity} \\
  & \lambda x. \lambda y. x && \text{constant} \\
  & \lambda x. \lambda y. x \text{ } y \text{ } y && \text{double application} \\
  & \lambda x. \lambda y. \lambda z. x (y \text{ } z) && \text{function composition}
\end{align*}

Having no built-in constant or primitive operators (e.g. numbers, arithmetic operations,
conditionals, loops, etc.), the only way to "compute" something is by function application, also
known as $\beta$-reduction and noted $\to_\beta$). It consists of replacing every instance of the
abstracted variable in the abstraction body by the effective argument. Following is an example of a
$\lambda$-term that is $\beta$-reduced twice:
\begin{align*}
  (\lambda x. (\lambda y. y) x) ((\lambda w. w) z)
    & \to_\beta (\lambda x. (\lambda y. y) x) z \\
    & \to_\beta (\lambda y. y) z \\
    & \to_\beta z
\end{align*}

This apparently simple operation hide a complex corner case: name clashes. Nothing prevent two
functions from using the same name for their abstracted variable. One can not simply replace every
variable with the same name when performing a substitution but must also take variable scopes into
account. Here is a simple example that demonstrate that a naïve approche can fail to preserve the
semantic of the original $\lambda$-term:

\begin{displaymath}
  (\lambda x. \lambda y. x) y \not\to_\beta \lambda y. y \\
\end{displaymath}

One solution to this problem is to rename function arguments, also known as $\alpha$-equivalence and
noted $=_\alpha$, prior to $\beta$-reduction. Here is a correct $\beta$-reduction for the previous
example:
\begin{align*}
  (\lambda x. \lambda y. x) y
    & =_\alpha (\lambda x. \lambda w. x) y \\
    & \to_\beta \lambda w. y \\
\end{align*}

The difference is important. Under the naïve approche, the $\lambda$-term was wrongly reduced to the
identity function while the correct reduction lead to a function returning the constant $y$.

Many constructions from high level programming languages can be encoded using only those basic
features.  Unary functions are already supported and n-ary functions can be straitforwardly emulate
by having a function return an other function, as was done in the previous examples:

\begin{displaymath}
  \text{n-ary function} \equiv \lambda x_1. \lambda x_2 \dots \lambda x_n. t
\end{displaymath}

An other common construction is a \emph{let binding} which serves to attach an identifier to a
complex expression. It can be emulated in $\lambda$-calculus with a single function abstraction:

\begin{displaymath}
  \text{let } x = y \text{ in } t \equiv (\lambda x. t) y
\end{displaymath}

Although significantly less obvious, it is also possible to express booleans only with functions.
The encoding is based on the idea that any use of booleans can be express with only three
primitives: a constant representing a true value, a constant representing false values and an
operation to choose between two alternatives:
\begin{align*}
  \text{true}
    & \equiv \lambda t. \lambda f. t \\
  \text{false}
    & \equiv \lambda t. \lambda f. f \\
  \text{if } b \text{ then } t \text{ else } e
    & \equiv \lambda b. \lambda t. \lambda f. b \text{ } t \text{ } f
\end{align*}

Using the rules already discussed, it is trivial to show that the following reduction is valid:

\begin{displaymath}
  \text{if true then } x \text{ else } y \to_\beta \dots \to_\beta x
\end{displaymath}

Other encodings exist for constructions such as numbers, list, datatypes, arbitrary recursion, etc.
For a more comprehensive introduction to the subject, \emph{An Introduction to Lambda Calculi for
Computer Scientists} \cite{hankin-2004-ILCCS} can serve as a good starting point.


\subsection{Type Systems}

Type systems are a syntactic method to prove the absence of certain behaviors. They differ from
testing in that they are exhaustive, automatic and compositional. In this context, exhaustiveness
means that each checked invarient is proved to hold for the complete program instead of focusing on
a single unit of code. Automaticness means that the process is non-interactive, i.e. the algorithm
does not need other informations than what it can find in the source code of the program.
Compositional means that proofs for individual components ca be use to satisfy an obligation about
the interaction of the components.

The kind of errors detected depends on the specific type system considered: it can range from
fairly simple to very complex. Examples of simple errors include typographical mistakes, usage of
values of the wrong kind and usage of undefined operations:

\begin{center}
  \begin{tabular}{m{3.5cm} | m{5.5cm}}
    $\text{add} : \mathbb{N} \to \mathbb{N} \to \mathbb{N}$ \newline
    $\text{true} : \mathbb{B}$ \newline
    add true true
    & Error in function application, "add" expects a $\mathbb{N}$ but a $\mathbb{B}$ was provided.
  \end{tabular}
\end{center}

With sufficiently powerfull type systems, specific requirements can be provide as type annotations.
Examples of such include that the output of a sorting function is a permutation of its input, that
the argument of an indexing operation on a list is in a valid range and that two multiplied matrixes
have compatible dimenssions. The more informations one put in the types, the more invalid programs
the type checker can catch, but the more work is required by the programmer to convince the
algorithm that the program fullfill it's specifications. An important design decision when defining
a type system is to find a tradeoff between those two conflicting properties.

Since they are often bundle with the compiler of a programming language and, thus, part of the
normal programming cycle, type systems allow early detection of programming errors. Moreover, the
diagnoses of type checkers can often pointed accuratly the source of the error, at the difference of
run-time tests where the effect of an error can sometime be visible only much farther in the code
when something start to go wrong.

An other important way in which type systems can be use is as an abstraction tool. It is widely
accepted that, in the context of large scale software developpement, programs should be split into
distinct modules that communicate through interfaces. Types are an natural fit to serve as such an
interface. Even in smaller scale programming, it is usefull to caracterise a datatype not by the
way it is implemented but by the different operations that can be perform on it. This focus on
operations led to the concept of abstract datatypes.

Types can also be used an invaluable maintenance tool. More informally, they serve as a checked
documentation of programs and, being simpler than the computation they caracterise, they can help
to reason about such computation on a higher level. But they can also serve a very practical
purpose, by checking which part of a program is affected by a change. If one decides to change the
arguments of a function or remove a field in a structure, a simple type checking pass will provide
an exhaustive list of the places that need to be updated to cope with the modification.

% Language safety is a very desirable property of programming language that can be achived by type
% systems.
%   * e.g. memory safe
%   * A safe language is completly defined by its programmer manual.
%     * C contains a lot of implementation defined and undefined behaviours.

Due to their static nature, type systems are conservatives in that they will always reject bad
programs at the expence of sometime rejecting good ones. A simple example of such limitation is the
following program that fails to typecheck, even if it happens that the complex boolean expression
will always evaluate to true at runtime:

\begin{displaymath}
  \text{if} <\text{complex-expression}> \text{then } 42 \text{ else } \text{true}
\end{displaymath}

Having only access to static information, a type checker only see that a boolean can take two
different values and, thus, must insure that the program is valide in both cases. In this example,
this implies that both branch of the \texttt{if} must be well-typed and that their type must be
compatible. It is the main goal of type theory scientists to develop type systems in which more more
valid programs are accepted while more invalid programs are rejected.

\subsection{Isabelle/HOL}

\input{Examples}

\section{Structure of the Formalization}
\label{sec:structure-of-formalization}

In this thesis, we formalize six chapters of the first two sections of \emph{TAPL}. Figure
\ref{fig:TAPL-toc} presents the table of contents of those two sections --- the formalized chapters
are in bold --- and Figure \ref{fig:TAPL-dependencies} presents the dependencies between the
chapters; a normal arrow implies a direct dependency while a dashed arrow only imply that the
knowledge learned in one chapter is reused.

\begin{figure}[h]
  \footnotesize
  \centering
  \begin{varwidth}{\textwidth}
    \begin{enumerate}[label=\Roman*]
      \itemsep 1pt
      \item Untyped Systems \hfill
        \begin{enumerate}[label=§ \arabic*]
          \setcounter{enumii}{2}
          \item \textbf{Untyped Arithmetic Expressions}
          \item An ML Implementation of Arithmetic Expressions
          \item \textbf{The Untyped Lambda-Calculus}
          \item \textbf{Nameless Representation of Terms}
          \item An ML Implementation of the Lambda-Calculus
        \end{enumerate}
      \item Simple Types \hfill
        \begin{enumerate}[label=§ \arabic*]
          \setcounter{enumii}{7}
          \item \textbf{Typed Arithmetic Expressions}
          \item \textbf{Simply Typed Lambda-Calculus}
          \item An ML implementation of Simple Types
          \item Simple Extensions
          \item Normalization
          \item References
          \item Exceptions
        \end{enumerate}
    \end{enumerate}
  \end{varwidth}
  \caption{Part I and II of \emph{TAPL}}
  \label{fig:TAPL-toc}
\end{figure}

\begin{figure}[h]
  \centering
  \begin{minipage}{0.49\textwidth}
    \centering
    \begin{tikzpicture}[>=triangle 45]
      \draw (0,0) rectangle (1,1) node[pos=0.5]{6};
      \draw (1.5,1.5) rectangle (2.5,2.5) node[pos=0.5]{5};
      \draw (3,0) rectangle (4,1) node[pos=0.5]{9};
      \draw (3,3) rectangle (4,4) node[pos=0.5]{3};
      \draw (4.5,1.5) rectangle (5.5,2.5) node[pos=0.5]{8};
      \draw [->] (1,1) -- (1.5,1.5); % 6 -> 5
      \draw [->] (3,1) -- (2.5,1.5); % 9 -> 5
      \draw [->, dashed] (2.5,2.5) -- (3,3); % 5 -> 3
      \draw [->] (4.5,2.5) -- (4,3); % 8 -> 3
      \draw [->, dashed] (4,1) -- (4.5,1.5); % 9 -> 8
    \end{tikzpicture}
    \captionof{figure}{Dependencies between the chapters of TAPL}
    \label{fig:TAPL-dependencies}
  \end{minipage}
  \begin{minipage}{0.49\textwidth}
    \centering
    \begin{tikzpicture}[>=triangle 45]
      \draw (0,3) rectangle (1,4) node[pos=0.5]{6};
      \draw (1.5,1.5) rectangle (2.5,2.5) node[pos=0.5]{5};
      \draw (3,0) rectangle (4,1) node[pos=0.5]{9};
      \draw (3,3) rectangle (4,4) node[pos=0.5]{3};
      \draw (4.5,1.5) rectangle (5.5,2.5) node[pos=0.5]{8};
      \draw [->] (1.5,2.5) -- (1,3); % 5 -> 6
      \draw [->, dotted] (3,1) -- (2.5,1.5); % 9 -> 5
      \draw [->, dashed] (2.5,2.5) -- (3,3); % 5 -> 3
      \draw [->] (4.5,2.5) -- (4,3); % 8 -> 3
      \draw [->, dashed] (4,1) -- (4.5,1.5); % 9 -> 8
    \end{tikzpicture}
    \captionof{figure}{Dependencies between the theory files}
    \label{fig:thys-dependencies}
  \end{minipage}
\end{figure}

The formalization closely follows the structure of the book. We provide one Isabelle theory file per
chapter and mainly introduce them in the same order. The only exception is the chapter on ``Nameless
Representation of Terms''. In the book, it is treated as a completely separate issue from the
chapters on $\lambda$-calculus (i.e. as an encoding that can be useful for an implementation). They
present it as a concrete alternative to the implicit $\alpha$-conversion they assume in their
proofs, which is nice for humans but not rigorous enough for a computer. Even though a formalization
is different from an implementation it has some similar requirements. Since we chose this nameless
representation for our formalization, we must diverge from the book and introduce this subject
before the untyped $\lambda$-calculus. Figure \ref{fig:thys-dependencies} presents the dependencies
between our Isabelle theory files.

It is possible to directly base the typed arithmetic expressions on untyped arithmetic expression
by importing the theory and reusing its definition. This reuse is possible because, for this
language, types are external to the representation of terms. The same argument applies to the
untyped $\lambda$-calculus. By contrast, for the typed $\lambda$-calculus, we need to alter the
representation of terms to add the typing annotation on abstraction variables, thus preventing the
reuse of the untyped $\lambda$-calculus theory. This is represented as a dotted arrow.


% sane default for proof documents
\parindent 0pt\parskip 0.5ex

%\part{Untyped Systems}
\input{Untyped_Arithmetic_Expressions}
\input{Nameless_Representation_Of_Terms}
\input{Untyped_Lambda_Calculus}

%\part{Typed Systems}
\input{Typed_Arithmetic_Expressions}
\input{Typed_Lambda_Calculus}

%\part{Conclusion}
\section{Conclusion}

In this thesis, we formalized a number of languages presented in \emph{Types and Programming
Languages} using the Isabelle/HOL interactive theorem prover. We started with a simple arithmetic
language of Booleans and natural numbers. We continued with the nameless representation of terms for
the $\lambda$-calculus, which we used as a basis for the pure untyped $\lambda$-calculus. For those
languages, we proved the determinacy of evaluation, the relation between values and normal form, the
uniqueness of normal form and the termination, or non-termination, of evaluation.

We then revisited both languages and augmented them with type systems. We proved the uniqueness of
types and the safety of the languages through the progress and preservation theorems. We also
demonstrated that the addition of types did not changed the semantics of the $\lambda$-calculus by
proving that types can be erased without affecting the evaluation of terms.

A formalization can be separated in three main elements: definitions, properties involving these
definitions and computer-checked proofs that these properties hold. In retrospective, expressing the
definitions and properties was the most important and difficult activity. Once the right
abstractions and the correct formulations for theorems were found, the proofs were usually fairly
simple: a good definition is worth three theorems! Conversely, a wrong abstraction or hypothesis
have led us to theorems very difficult to prove, or even to properties that ended up not being
theorems at all.

In this report, we focused our attention on the definitions and theorems, highlighting the
differences with the book. The complete Isabelle/HOL theories provided along with this
report\footnote{\url{https://github.com/authchir/log792-type-systems-formalization}} contain more
examples, exercises and less important theorems.


\nocite{*}
\bibliographystyle{alpha}
\bibliography{root}

% optional bibliography
%\bibliographystyle{abbrv}
%\bibliography{root}

\end{document}

%%% Local Variables:
%%% mode: latex
%%% TeX-master: t
%%% End:
